\section{Code de mise en \oe{}uvre de l'estimateur}
\label{ann:B}

Cette annexe présente les principales routines Stata utilisées pour l'estimation
PSM-DD. Le code complet est organisé en scripts numérotés dans le répertoire
\texttt{code/stata/} du projet (un script unique \texttt{tout.do} regroupe
l'ensemble du pipeline)~;
les erreurs-types sont clusterisées au niveau de la grappe.

\subsection{Statut de traitement}

\begin{verbatim}
/* 04_panel.do - beneficiaires 2018 vs non beneficiaires */
use "$TEMP/traitement_2018.dta", clear
rename D D_2018
merge 1:1 grappe menage using "$TEMP/traitement_2021.dta", ///
    keepusing(D) keep(match) nogenerate
rename D D_2021
gen byte D_stable = .
replace D_stable = 1 if D_2018 == 0 & D_2021 == 1
replace D_stable = 0 if D_2018 == 0 & D_2021 == 0
/* menages deja beneficiaires en 2018 (D_stable == .)
   exclus du panel causal */
\end{verbatim}

\subsection{Construction du panel d'enfants suivis}

\begin{verbatim}
/* Cle individuelle entre vagues : identifiant 2018 precharge en 2021.
   Le rang dans le roster 2021 ne designe pas le meme individu qu'en 2018. */
use "$BASE_2021/s01_me_sen2021.dta", clear
keep grappe menage membres__id s01qpreload_pid
rename membres__id     numind
rename s01qpreload_pid numind_2018
keep if !missing(numind_2018)
save "$TEMP/lien_individus.dta", replace

/* Jointure des deux vagues sur la cle individuelle */
use "$TEMP/enfants_dep_2021.dta", clear
merge m:1 grappe menage numind using "$TEMP/lien_individus.dta", ///
    keepusing(numind_2018) keep(match) nogenerate
drop numind
merge 1:1 grappe menage numind_2018 using `enf18_psm', keep(match) nogenerate
egen long enfid = group(grappe menage numind_2018)
\end{verbatim}

\subsection{Estimation du score de propension au niveau enfant}

\begin{verbatim}
/* Logit ENFANT sur t=0 : probabilite qu'un enfant vive dans un
   menage beneficiaire. Covariables du menage + sexe et age de l'enfant. */
logit D c.hhsize c.log_pcexp i.milieu i.region ///
         c.hgender c.hage i.heduc i.hmstat i.hcsp ///
         i.sexe18 c.age18, vce(cluster grappe) nolog
predict pscore, pr
\end{verbatim}

\subsection{Appariement et équilibre des covariables}

\begin{verbatim}
/* Appariement k plus proches voisins (k=4, avec remise), niveau enfant */
psmatch2 D, pscore(pscore) neighbor(4) common
pstest hhsize log_pcexp i.milieu i.region hgender hage ///
       i.heduc i.hmstat i.hcsp i.sexe18 age18, both
rename _weight weight_knn
keep enfid grappe menage numind_2018 D pscore weight_knn _support
save "$TEMP/pscore_knn.dta", replace
\end{verbatim}

\subsection{Combinaison appariement et double différence}

\begin{verbatim}
/* Panel d'enfants en format long : chaque enfant porte SON poids
   d'appariement a ses deux dates. */
use "$TEMP/pscore_t0.dta", clear
merge 1:1 enfid using "$TEMP/pscore_knn.dta", ///
    keepusing(weight_knn _support) nogenerate
reshape long pauvre_MODA nb_dep groupe_moda sexe age ///
             dim_assai dim_eau dim_logem dim_nutri ///
             dim_sante dim_protect dim_educ, i(enfid) j(periode)
gen byte t = (periode == 21)
keep if !missing(weight_knn) & weight_knn > 0

/* ATT PSM-DD : coefficient de l'interaction t x D */
regress pauvre_MODA i.t##i.D [aw=weight_knn], vce(cluster grappe)
lincom 1.t#1.D
\end{verbatim}

\subsection{Effets par dimension}

\begin{verbatim}
/* 07_effets_dim.do */
local dims assai eau logem nutri sante protect educ
foreach dim of local dims {
    quietly regress dim_`dim' i.t##i.D [aw=weight_knn], ///
        vce(cluster grappe)
    quietly lincom 1.t#1.D
    di "dim_`dim' : ATT=" %8.4f r(estimate) ///
       "  p=" %6.4f r(p)
}
\end{verbatim}

Le script \texttt{tout.do} regroupe l'ensemble du pipeline en sections
successives~: exploration des bases (\texttt{01\_visitation}), variable de
traitement (\texttt{02\_traitement}), indicateurs de pauvreté
(\texttt{03\_deprivation}), construction du panel (\texttt{04\_panel}),
estimation PSM-DD au niveau enfant (\texttt{05\_psm\_dd}), statistiques
descriptives (\texttt{06\_stats\_desc}), effets par dimension
(\texttt{07\_effets\_dim}) et tests de validité
(\texttt{09\_placebo\_attrition}). Toutes les valeurs reportées dans le
mémoire proviennent du journal d'exécution de ce script.
